% Box transpose — CS 8045 project report. Build with `make` — see README.md.
%
% Document skeleton only. Everything here is either a package/style decision
% (style/) or the order sections appear in (sections/) — no prose, no floats,
% no measured numbers live in this file.
%
% This is the results paper, not the proposal. The proposal is a separate
% deliverable (presentations/), and the milestone and simulation-setting
% sections it needed are gone from here.
%
% Section order:
%   sections/0-title        team and responsibilities
%   sections/1-problem      what .contiguous() costs and why it matters
%   sections/method/        box swaps, fixed points, the search
%   sections/evaluation/    setup, the two tables, the three regimes
%   sections/3-limitations  where the method is the wrong tool
%   sections/4-related      prior work, and what we do not claim
%   sections/5-conclusion
%
% NOTE ON PATHS: LaTeX resolves \input relative to the directory pdflatex runs
% in (paper/), not relative to the file doing the including. Every path here,
% including inside nested section files, is written from that root.
\documentclass[letterpaper,twocolumn,10pt]{article}

% House preamble, carried over unchanged from the smartcontiguous paper so this
% document behaves identically to that one. Each file carries the reasoning for
% its own settings; read those comments before changing anything here.
% The style file loads xcolor itself, so the table option has to be handed over
% before that happens. It brings in colortbl, for row and cell backgrounds.
% Table 4 used it for a reversed-out legend band until 2026-08-23; nothing in
% the paper draws on a colour now, so this can go if nothing new wants it.
\PassOptionsToPackage{table}{xcolor}
\usepackage{style/usenix-2020-09}

\usepackage{tikz}
\usepackage{helvet}
% arrows.meta is loaded alongside arrows, not instead of it: the older figures
% keep their >=stealth' tips, while fig-thread-race sizes its heads in absolute
% pt. Old style tips scale with the line width, which made the 3pt connectors
% in that figure read as blobs rather than arrows.
\usetikzlibrary{arrows,arrows.meta,backgrounds,decorations.pathreplacing,patterns.meta}
\usepackage{amsmath}
\usepackage{amssymb}
\usepackage{bm}
\usepackage{graphicx}
\usepackage{booktabs}
\usepackage{array}
\usepackage{algorithm}
\usepackage{algpseudocode}
\usepackage{float}
\usepackage{makecell}
\usepackage{comment}

% The bibliography carries long URLs and \texttt{} commit hashes. In a narrow
% two-column measure those are unbreakable boxes, so TeX pads the inter-word
% spaces around them instead, which shows up as rivers of white in an entry.
% xurl lets a URL break at any character; emergencystretch gives the paragraph
% builder a last-resort pass so it stops stretching good lines to save a bad one.
\usepackage{xurl}

% breaklinks, because the usenix style loads hyperref with no options and the
% pdftex default is false. A link that breaks across lines then gets ONE
% bounding box spanning everything between its start and its end, and xurl
% breaks long URLs freely, so a citation whose URL straddled a column break
% produced a clickable rectangle 200pt tall covering the whole of Figure 5.
% With breaklinks the link is emitted one rectangle per line.
\hypersetup{breaklinks=true}

% Every column is set to the full 9in measure, the way the FAST papers in
% role-models/ are.  \raggedbottom left each column ending wherever its last
% line fell, which reads as a ragged, uneven foot across a spread.  The floats
% are placed so little is left to absorb, so the stretch stays invisible.
\flushbottom
\emergencystretch=2em

% A page holding only floats ignores [t]: LaTeX centres its contents between
% \@fptop and \@fpbot, which both default to 0pt plus 1fil. That left the
% appendix table stranded in the middle of an otherwise blank page. Dropping the
% stretch at the top makes a float page start at the top margin, which is what
% [t] already asks for everywhere else.
\makeatletter
\setlength{\@fptop}{0pt}
\setlength{\@dblfptop}{0pt}   % table* spans both columns and uses its own glue
\makeatother

% Number and display two levels below subsubsection, so the results section can
% show its own structure (breadth vs depth) instead of relying on prose to
% signal it. This also promotes the eight pre-existing \paragraph labels in the
% design section from run-in italicised asides to numbered headings, which is
% the same improvement applied consistently rather than only where it was asked
% for. Their trailing periods were stripped to match.
\setcounter{secnumdepth}{5}
\makeatletter
\renewcommand\paragraph{\@startsection{paragraph}{4}{\z@}%
  {2.0ex \@plus1ex \@minus.2ex}{0.6ex \@plus.2ex}%
  {\normalfont\normalsize\bfseries}}
\renewcommand\subparagraph{\@startsection{subparagraph}{5}{\z@}%
  {1.8ex \@plus1ex \@minus.2ex}{0.5ex \@plus.2ex}%
  {\normalfont\small\bfseries}}
\makeatother

% Caption convention borrowed from the FAST'26 papers in role-models/: the
% opening sentence is set bold as a headline, the rest reads as ordinary prose.
% It gives every float a scannable one-line summary. Defined once so the whole
% paper can be switched back by editing this line alone.
\newcommand{\capheadline}[1]{\textbf{#1}}

% Visible, greppable placeholder. Deliberately loud so an unfinished section
% cannot be mistaken for finished prose in a draft read. Delete the macro and
% every \placeholder call before submission; `grep -n placeholder` finds them.
\newcommand{\placeholder}[1]{%
  \par\medskip\noindent\fbox{\parbox{0.95\linewidth}{%
    \footnotesize\textbf{[PLACEHOLDER]} #1}}\par\medskip}

\newcommand{\pname}[1]{\textsc{SMC}{#1}}
\newcommand{\aname}[1]{\textsc{CS}{#1}}

% Breakable pseudocode listings for the appendix. An algorithm float cannot
% split across a column, and each appendix listing is most of a column tall, so
% floating them left a large hole wherever one did not fit. These two macros
% reproduce the ruled float's appearance while letting the body break across
% columns like ordinary text. \refstepcounter keeps "Algorithm N" numbering and
% \ref working from the body.
\makeatletter
\newcommand{\algstart}[2]{%
  \par\addvspace{1.4ex}%
  \refstepcounter{algorithm}%
  \noindent\rule{\linewidth}{0.9pt}\par\nobreak\vspace{0.5ex}%
  {\interlinepenalty=10000\noindent{\normalsize\textbf{Algorithm \thealgorithm:} #1}\label{#2}\par}%
  \nobreak\vspace{0.5ex}%
  \noindent\rule{\linewidth}{0.4pt}\par\nobreak\vspace{0.4ex}}
\newcommand{\algfinish}{%
  \par\nobreak\vspace{0.4ex}%
  \noindent\rule{\linewidth}{0.9pt}\par\addvspace{1.4ex}}
\makeatother


% Generated number macros, from the producers under pipeline/: the prose
% claims (quantities/), and whatever each table and figure renderer computed
% for its own caption. Loaded here rather than inside the section that cites
% them, so any section can quote the same live values. Add one \input line per
% _generated/*.defs.tex as you add generators; pipeline/quantities/audit.py
% fails the build if a hand-written section file carries a literal none of
% them accounts for.
% GENERATED by pipeline/quantities/gen_model_defs.py -- do not edit by hand.
% Regenerate after any change to src/_plan.py or raw_data/dpost_sweep/.
\newcommand{\modelMinMoved}{0.50}
\newcommand{\modelMinMovedRows}{2}
\newcommand{\modelMinMovedCols}{2}
\newcommand{\modelErrIn}{2.4}
\newcommand{\modelErrOut}{64}
\newcommand{\modelMaxSteps}{6}
%
% GENERATED by pipeline/quantities/gen_results_defs.py -- do not edit by hand.
% Regenerate after any change to raw_data/dpost_sweep/ or raw_data/hard_perms/.
\newcommand{\resCases}{87}
\newcommand{\resRankLo}{3}
\newcommand{\resRankHi}{10}
\newcommand{\resElements}{268{,}435{,}456}
\newcommand{\resTensorMB}{1{,}024}
\newcommand{\resDevice}{NVIDIA GeForce RTX 2070}
\newcommand{\resBoxMs}{10.21}
\newcommand{\resContMs}{7.70}
\newcommand{\resIttpdMs}{80.49}
\newcommand{\resGlMs}{82.91}
\newcommand{\resGvMs}{97.39}
\newcommand{\resBoxAux}{144}
\newcommand{\resGvAuxMB}{48}
\newcommand{\resWinIttpd}{87}
\newcommand{\resWinIttpdOf}{87}
\newcommand{\resWinGl}{87}
\newcommand{\resWinGlOf}{87}
\newcommand{\resWinGv}{79}
\newcommand{\resWinGvOf}{87}
\newcommand{\gateCases}{21}
\newcommand{\gateWins}{21}
\newcommand{\gateShare}{24}
\newcommand{\gateMedianX}{1.73}
\newcommand{\gateBestX}{2.05}
\newcommand{\gateMinDpost}{1{,}024}
\newcommand{\gateMaxSide}{4}
\newcommand{\hardCases}{15}
\newcommand{\hardBigCases}{4}
\newcommand{\hardBigWins}{4}
\newcommand{\hardSmallCases}{11}
\newcommand{\hardSmallWins}{0}
\newcommand{\hardSmallBoxLo}{146}
\newcommand{\hardSmallBoxHi}{314}
\newcommand{\hardSmallContLo}{7}
\newcommand{\hardSmallContHi}{14}
\newcommand{\hardSmallSlowdown}{30}
%
% GENERATED by pipeline/tables/gen_table_methods.py -- do not edit by hand.
% Regenerate after any change to raw_data/dpost_sweep/.
\newcommand{\methTabCases}{87}
\newcommand{\methTabIttpdX}{7.9}
\newcommand{\methTabGvX}{9.5}
\newcommand{\methTabGlX}{8.1}
%
% GENERATED by pipeline/tables/gen_table_tiers.py -- do not edit by hand.
% Regenerate after any change to raw_data/dpost_sweep/ or raw_data/hard_perms/.
\newcommand{\tierBeatsCases}{22}
\newcommand{\tierBeatsX}{1.73}
\newcommand{\tierInplaceCases}{69}
\newcommand{\tierSmallCases}{11}
\newcommand{\tierSmallX}{29}
%

\begin{document}

\title{\Large \bf Box Transpose: In-Place N-Dimensional Tensor Permutation on GPUs}

% One author block, not two. \and would set the course line as a second
% author, which is what it looked like. Affiliation and the rest of the team
% are in sections/0-title, where a \placeholder can be a block without
% stretching the title.
\author{
{\rm Kaveen Jayamanna \quad Alexis Newman}\\
Kennesaw State University\\
{\rm CS 8045 Advanced Algorithms}
}

\maketitle

\begin{abstract}
Changing the order of a tensor's axes in PyTorch is free. Fixing the layout
afterwards is not. \texttt{.contiguous()} makes a full copy, so peak memory
doubles. On a GPU that is almost full, the job fails there.

We present \emph{box transpose}. It fixes the layout in place. A shortest-path
search breaks the permutation into swaps of neighbouring groups of axes. Each
swap is a batched matrix transpose, which we carry out by following its cycles.
The only extra memory is the list of where each cycle starts.

A copy has to write every value, even the ones already in the right place. We
skip them. So box transpose can move less than the whole tensor, and a copy
never can. How much it moves follows from the shape in one line of arithmetic.
That predicts the speedup over \texttt{.contiguous()} before any data is
touched, with no constant measured on a GPU. The prediction lands within about
\modelErrIn\% of the measured time.

We tested \resCases{} randomly drawn shapes at ranks \resRankLo{} to
\resRankHi{}, each \resTensorMB{}\,MB, on a \resDevice{}. Box transpose needed
\resBoxAux{} bytes of extra memory in the median case. Against three published
in-place methods it was faster on \resWinIttpd{}, \resWinGl{} and \resWinGv{}
of \resCases{} cases, by \methTabIttpdX{}$\times$ to \methTabGvX{}$\times$ in
the median. When the permutation leaves a large group of axes at the back
alone, a check on the shape picks out the cases where we beat the copy itself.
We won all \gateCases{} of them, \gateMedianX{}$\times$ faster in the median.
When that group is too small, we are about \hardSmallSlowdown{}$\times$ slower
than the copy, and the same check says so in advance.
\end{abstract}

\section{The Problem}
\label{sec:problem}

\subsection{Changing the order of axes is free}

Deep learning code changes the order of a tensor's dimensions often. In
PyTorch, \texttt{permute} and \texttt{transpose} do
this~\cite{pytorch_permute_docs, pytorch_transpose_docs}. They are free. They
change a label. They do not move any data.

\subsection{Fixing the layout is not free}

After that, the data is in the wrong order in memory. Some operations will not
run on it. Others run, but they use a slow path. They do not warn you.

The normal fix is \texttt{.contiguous()}~\cite{pytorch_contiguous_2016}. It
makes a copy. It asks for a new tensor. It writes every value into the right
place. Then it throws the old tensor away.

\subsection{The copy is the problem}

While the copy runs, both tensors are in memory. So a tensor of $B$ bytes needs
at least $2B$ bytes. The old one must stay until it has been read. The new one
must exist to be written into. This is what a copy means. You cannot avoid it
by changing the loop.

If the machine is almost full, the job fails here. A 2026 PyTorch design
document says non-contiguous tensors now appear all over modern
training~\cite{pytorch_rfc_comms_2026}. It says some communication routines
will not take them. It says the usual fixes either make a full copy or need a
buffer as big as the data.

There is a second problem. The copy is a large block of memory. It only lives
for a short time. It appears in the middle of work that already uses a lot of
memory. Earlier work shows this pattern causes
fragmentation~\cite{rajbhandari2020zero}. Memory is free, but it is in small
pieces. The next request needs one big piece and fails.

\subsection{What we do}

We fix the layout in place. We never ask for a second tensor. We break the
reordering into small swaps. Each swap exchanges two groups of axes that sit
next to each other. We store only a few hundred bytes while we work.

Doing this in place is not new. Section~\ref{sec:related} lists the earlier
work. Two things here are new. In more than two dimensions, this can do
\emph{less work than a copy}. And we can tell which cases those are by looking
at the shape, before we touch any data.

\section{How It Works}
\label{sec:method}

\subsection{Swapping boxes}
\label{sec:boxes}

A \emph{box} is one or more axes that sit next to each other. We treat a box as
one unit. Their sizes multiply together, so a box acts like a single axis of
that size.

A \emph{swap} exchanges two boxes that sit next to each other. Any reordering
of axes can be done as a list of swaps.

Here is why that helps. Cut the axis list into four parts: a front part, a left
box, a right box, and a back part. Call their sizes $D_{pre}$,
$\mathit{rows}$, $\mathit{cols}$ and $D_{post}$. Swapping the two middle boxes
is the same as transposing a $\mathit{rows} \times \mathit{cols}$ matrix. We do
it once for each of the $D_{pre}$ front positions. Each cell of that matrix is
a run of $D_{post}$ values that sit next to each other.

This idea is not ours. It is the $\rho(i,j,k)$ operation from
ITTPD~\cite{cheng2025ittpd}. Section~\ref{sec:related} says more. What is ours
is how we run a swap, and how we choose which swaps to make.

\subsection{Doing a swap with no spare room}

You cannot move a value straight to where it belongs. Something else is already
there. So you pick a value up. You put it where it goes. Now you are holding
whatever was in that spot. You keep going. In the end you come back to where
you started.

That loop is called a \emph{cycle}. Walking it is called cycle
following~\cite{gustavson2012cycleleader,dummit2003abstract}. The only thing
you store is the one value in your hand. On a GPU that is one register per
thread.

In two dimensions this does not spread well over a GPU. You can only run as
many cycles at once as there are cycles.
Catanzaro~et~al.~\cite{catanzaro2014decomposition} say this plainly. A shape
with nine cycles gets nine-way parallelism, even on a machine with hundreds of
processors.

With more dimensions that limit goes away. We repeat the swap $D_{pre}$ times.
Each value we move is $D_{post}$ values wide. So the work splits into
$D_{pre} \times \mathit{cycles} \times D_{post}$ pieces that can run at the
same time. Those two extra numbers come from the axes around the swap. A
two-dimensional method does not have them.

\subsection{Skipping cells that are already correct}
\label{sec:fixedpoints}

Some cells are already in the right place. A transpose never moves the diagonal
of a square matrix.

\textbf{We skip those cells. A copy cannot.} \texttt{.contiguous()} is building
a new tensor. It must write every position, even the ones that were already
correct.

There is a formula for how many cells stay put. The transpose sends position
$p$ to $(p \bmod \mathit{cols})\,\mathit{rows} + \lfloor p/\mathit{cols}\rfloor$.
In the middle of the array this is the same as multiplying by $\mathit{cols}$
modulo $M = \mathit{rows}\cdot\mathit{cols}-1$. A cell stays put when
$p(\mathit{cols}-1) \equiv 0 \pmod M$. That has
$\gcd(\mathit{rows}-1,\mathit{cols}-1)$ answers. The last position also stays
put. So the share of the tensor a swap touches is

\begin{equation}
\mathit{moved} = 1 - \frac{\gcd(\mathit{rows}-1,\,\mathit{cols}-1) + 1}
                         {\mathit{rows}\cdot\mathit{cols}} .
\label{eq:moved}
\end{equation}

That is our whole cost model. A copy reads every value and writes every value.
It always moves $2N$ values' worth of traffic. In the cases where our method is
worth using, both methods move data at the same speed. So the speed cancels
out. Only the amount is left. That makes Equation~(\ref{eq:moved}) our
predicted time as a share of the copy's time. No number in it was measured on a
GPU. If it is below one, we do less work than the copy.

Now the main point. Any method that builds output positions must go over the
tensor at least once. It has to write every position. Ours does not. Swapping a
$\modelMinMovedRows \times \modelMinMovedCols$ block moves $\modelMinMoved$ of
the tensor. No single swap can do better. The reason is that
$\gcd(\mathit{rows}-1,\mathit{cols}-1)$ is at most
$\min(\mathit{rows},\mathit{cols})-1$, so the skipped share is at most
$1/\max(\mathit{rows},\mathit{cols})$.

Two things follow. Only a method that skips can do part of a pass. And two
swaps can never add up to less than one pass. So \emph{beating
\texttt{.contiguous()} and using a single swap are the same thing}.

\subsection{Choosing which swaps to make}
\label{sec:planning}

There are many lists of swaps that reach the same order. They do not cost the
same.

We build a graph. Each way of ordering the axes is a point. Each legal swap is
a line between two points. The weight on that line is how many bytes the swap
moves. We know that number exactly from Equation~(\ref{eq:moved}). Then we run
Dijkstra's algorithm to find the cheapest route.

We do not remove any lines. We do not apply any cutoff to them. We only reject
a finished route. If the cheapest route needs more than $\modelMaxSteps$ swaps,
we return nothing. The caller then uses \texttt{.contiguous()}. Every swap
rewrites the whole buffer, so a long route moves several tensors' worth of data.

The search does not make us faster than a copy. Section~\ref{sec:fixedpoints}
shows that needs a single swap, and you do not need a search to find one. The
search matters for the longer routes. Nothing beats a copy there, but the
search keeps us ahead of the other in-place methods.

\subsection{Nothing is tuned}

An older version of this work used three numbers measured on one GPU. There was
a throughput knee, a routing cutoff, and an average bandwidth. All three are
gone. The bandwidth cancelled out. The knee became a question about the shape
instead of a measurement.

We can check that the model knows where it works. Inside the region we claim,
Equation~(\ref{eq:moved}) predicts the measured ratio to within
$\modelErrIn\%$. Outside it, the error is about $\modelErrOut\%$. That is why
the region exists.

\section{Results}
\label{sec:eval}

\subsection{How we tested}
\label{sec:setup}

We did not choose shapes by hand. Take a tensor with
$N = \prod_{i=1}^{m} p_i^{a_i}$ values and rank $r$. Every valid shape comes
from splitting each prime's exponent across the $r$ dimensions, so that
$\sum_{j=1}^{r} b_{ij} = a_i$:

\begin{equation}
s = \left( \prod_{i=1}^{m} p_i^{b_{i1}},\ \prod_{i=1}^{m} p_i^{b_{i2}},\
           \ldots,\ \prod_{i=1}^{m} p_i^{b_{ir}} \right).
\label{eq:shapes}
\end{equation}

We draw from that whole set. So nobody picked the cases. We drop shapes with an
axis of size one. Such an axis can move anywhere without touching a byte.

The main test has \resCases{} cases at ranks \resRankLo{} to \resRankHi{}.
Every tensor has \resElements{} float32 values, which is \resTensorMB{}\,MB. We
ran them on a \resDevice{}. Each shape is paired with a permutation that leaves
a group of axes at the back alone. That is what our method is built for.

We ran a second test with \hardCases{} permutations picked at random. That
covers what our method is not built for. Section~\ref{sec:limits} reports it.

We compare against PyTorch \texttt{.contiguous()}. We also compare against our
rebuilds of the three published in-place methods.
Section~\ref{sec:related} lists them. We time every method with CUDA events
after one warmup run. We check every answer against
\texttt{permute(*perm).contiguous()} before it counts.

Our method finished and gave the right answer on every case in both tests. It
turned nothing away.

% Table (tab:methods). Body and caption macros: _generated/methods-table.tex
% and .defs.tex, written by pipeline/tables/gen_table_methods.py from
% raw_data/dpost_sweep/.
%
% table*, not table. Five columns -- a method name plus two byte-counts wide
% enough to carry a unit -- have a natural width well past the 241pt of a
% single USENIX column. tabular* does NOT shrink to fit: it only distributes
% slack, so when the content is already too wide it runs past the column and
% over the neighbouring text WITHOUT raising an Overfull warning. That is why
% this has to be caught by eye, and why the full measure is the fix rather
% than \small-ing the type further.
\begin{table*}[!t]
\centering
\caption{\capheadline{Every method, every case, no routing.}
\methTabCases{} shape and permutation pairs at ranks \resRankLo{}--\resRankHi{},
\resTensorMB{}\,MB per tensor, on a \resDevice{}. We run our method on all of
them, including the ones our own rule would turn away.}
\label{tab:methods}
\small
\begin{minipage}{0.72\linewidth}
\centering
% GENERATED by pipeline/tables/gen_table_methods.py -- do not edit by hand.
% Regenerate after any change to raw_data/dpost_sweep/.
\begin{tabular}{lrrrr}
\toprule
Method & Median & Median aux & Peak aux & Box wins \\
\midrule
PyTorch \texttt{.contiguous()} & 7.70 & 1,024\,MB & 1,024\,MB & 32/87 \\
Box transpose (ours) & \textbf{10.21} & \textbf{144\,B} & \textbf{10.6\,KB} & -- \\
ITTPD~\cite{cheng2025ittpd} & 80.49 & 0 & 0 & 87/87 \\
Gomez-Luna~\cite{gomezluna2016inplace} & 82.91 & 16\,B & 512\,B & 87/87 \\
Gustavson~\cite{gustavson2012cycleleader} & 97.39 & 48\,MB & 128\,MB & 79/87 \\
\bottomrule
\end{tabular}

\vspace{1.5pt}
\parbox{\linewidth}{\fontsize{6}{6.8}\selectfont\raggedright Middle-of-the-range values over every case, with no routing: we run our method on all of them, including the ones our own rule would turn away. \textbf{Box wins} counts cases where box transpose is faster than that method. Against \texttt{.contiguous()} that count is low here precisely because no routing is applied; Section~\ref{sec:gate} reports it under the gate. ITTPD's memory figure comes from its own planner, not from watching the allocator. On these shapes every block fits in fast on-chip memory, so it never asks for anything global.}
%
\end{minipage}
\end{table*}


\subsection{Against everyone}

Table~\ref{tab:methods} shows all \methTabCases{} cases. We used no routing
here. We ran our method even on the cases our own rule says to skip.

In the middle of the range we are \methTabIttpdX{}$\times$ faster than ITTPD,
\methTabGlX{}$\times$ faster than Gomez-Luna, and \methTabGvX{}$\times$ faster
than Gustavson. Case by case, we are faster on \resWinIttpd{} of
\resWinIttpdOf{}, on \resWinGl{} of \resWinGlOf{}, and on \resWinGv{} of
\resWinGvOf{}.

The memory column matters too. \texttt{.contiguous()} needs a second tensor
every time. We need \resBoxAux{} bytes. Gustavson needs \resGvAuxMB{}\,MB.
ITTPD needs none. That is real, not a mistake. On these shapes its blocks fit
in fast on-chip memory, so it never asks for global memory. Against ITTPD we
win on speed, not on memory. We do not claim otherwise.

% Table (tab:tiers). Body and caption macros: _generated/tiers-table.tex and
% .defs.tex, written by pipeline/tables/gen_table_tiers.py from both sweeps.
%
% table*, not table: seven columns of data do not fit a 241pt column.
\begin{table*}[!t]
\centering
\caption{\capheadline{Three cases, and who wins each.}
Both sweeps together, split by the two numbers from
Section~\ref{sec:method}. $\tau$ is how long a run has to be before the memory
system reads it efficiently, \gateMinDpost{} values. The first row is where we
beat the copy. The last is where our method is the wrong tool.}
\label{tab:tiers}
\small
% GENERATED by pipeline/tables/gen_table_tiers.py -- do not edit by hand.
% Regenerate after any change to raw_data/dpost_sweep/ or raw_data/hard_perms/.
\begin{tabular*}{\linewidth}{@{\extracolsep{\fill}}lrrrrrr@{}}
\toprule
 & & \multicolumn{3}{c}{Median ms} & \multicolumn{2}{c}{Box faster than} \\
\cmidrule(lr){3-5}\cmidrule(lr){6-7}
Regime & Cases & Box & Copy & Best other & Copy & Best other \\
\midrule
$D_{post}\!\ge\!\tau$, one small square swap & 22 & 4.03 & 7.68 & 8.16 & 22/22 & 22/22 \\
$D_{post}\!\ge\!\tau$, otherwise & 69 & 12.48 & 7.69 & 22.37 & 11/69 & 61/69 \\
$D_{post}\!<\!\tau$ & 11 & 228.12 & 7.48 & 74.28 & 0/11 & 0/11 \\
\bottomrule
\end{tabular*}

\vspace{1.5pt}
\parbox{\linewidth}{\fontsize{6}{6.8}\selectfont\raggedright \textbf{Copy} is PyTorch \texttt{.contiguous()}. \textbf{Best other} is whichever in-place method was fastest on that case. Both tests are worked out from the shape before any data is touched. The last row is where our method falls down: a permutation that moves the innermost axis forces a swap with $D_{post}=1$, each value we move is then a lone number, and the published in-place methods are faster. We report this rather than route around it.}
%
\end{table*}


\subsection{Three cases}
\label{sec:gate}

Table~\ref{tab:tiers} splits both tests by the two numbers from
Section~\ref{sec:method}. We work out both from the shape before touching any
data.

The back group is large and the swap is one small square block. Here we are
faster than everyone, including the copy. We won \tierBeatsCases{} of
\tierBeatsCases{} cases, at \tierBeatsX{}$\times$ in the middle of the range.
In the main test alone, our rule fires on \gateShare{}\% of shapes and wins
\gateWins{} of \gateCases{}. That is \gateMedianX{}$\times$ in the middle and
\gateBestX{}$\times$ at best.

The back group is large but the swap does not qualify. Here we lose to the
copy. We are still the best in-place choice, on \tierInplaceCases{} cases.

The back group is too small. Here we are the worst. Section~\ref{sec:limits}
covers it.

\subsection{The prediction was right}

We did not fit the boundary to the results. Equation~(\ref{eq:moved}) says a
swap of a $\modelMinMovedRows \times \modelMinMovedCols$ block moves
$\modelMinMoved$ of the tensor. So it should take $\modelMinMoved$ of the
copy's time. From rank \resRankLo{} to rank \resRankHi{}, the measured times
land on that value. The error inside the region is about $\modelErrIn\%$.

The speedup comes from arithmetic, not from tuning. It does not depend on how
well we wrote the kernel. It does not depend on which GPU we used. It depends
only on how much of the tensor was already in the right place.

\section{Where It Falls Down}
\label{sec:limits}

\subsection{A small back group means no speed}

Some permutations move the innermost axis. Then every way of breaking the job
up has a swap with nothing behind it. For that swap, $D_{post}=1$.

Now each value we move is a single number on its own. The memory system reads a
whole cache line and uses four bytes of it. There is no way to group the work.

Our second test measures this. We used \hardCases{} permutations picked at
random. In \hardBigCases{} of them a back group survived. We were the best
in-place method on \hardBigWins{} of those \hardBigCases{}. In the other
\hardSmallCases{} no back group survived. We were the worst in-place method on
every one. We took \hardSmallBoxLo{} to \hardSmallBoxHi{}\,ms. The copy took
\hardSmallContLo{} to \hardSmallContHi{}\,ms. We were about
$\hardSmallSlowdown\times$ slower.

We still finish. We still ask for no extra memory. But here the published
in-place methods are several times faster than we are. A real system would send
these cases to one of them. We have not built that. We report the gap instead
of hiding it.

Rank~2 is always in this group. A plain matrix transpose has no back axes to
leave alone.

\subsection{More than one swap loses to the copy}

Section~\ref{sec:fixedpoints} shows a single swap always moves at least
$\modelMinMoved$ of the tensor. So a route with two or more swaps moves at
least as much as a copy does. We are often still the best in-place choice
there. But we are not faster than copying. A better search would not change
that.

\subsection{Planning is slow at high rank}

The search runs over orderings of the axes. It stops as soon as it reaches the
target. So the common cases, which need one swap, come back at once. But a hard
permutation at rank~8 took several seconds to plan.

A route depends only on the shape and the permutation. So we work it out once
and reuse it. A much higher rank would need a search that aims at the target
instead of spreading out in every direction.

\subsection{One GPU}

Every time in this paper comes from one \resDevice{}. The main speedup is
arithmetic. Equation~(\ref{eq:moved}) has no hardware term in it. So it should
hold on other machines. The comparisons against other systems are
measurements, and those need running again elsewhere.

We should also say this. Our rebuild of ITTPD runs its on-chip path about ten
times slower than a plain copy on these shapes. That may be our rebuild rather
than the published method. If so, our lead over it is smaller than we report.

\section{Earlier Work}
\label{sec:related}

\subsection{Transposing a matrix in place}

Catanzaro, Keller and Garland~\cite{catanzaro2014decomposition} split a matrix
transpose into row and column shuffles, with a rotation in between. Their
method needs a temporary buffer. Its cycle-walking step runs one block per
cycle. So it can only spread as wide as the number of cycles. They say this
themselves.

Gustavson, Karlsson and K\aa{}gstr\"{o}m~\cite{gustavson2012cycleleader} work
out where each cycle starts using number theory. They do not have to search for
it. We use the same idea for square blocks. There a transpose is its own
inverse, so the test becomes one comparison.

Sung et al.~\cite{sung2012pttwac} and G\'{o}mez-Luna et
al.~\cite{gomezluna2016inplace} pad dimensions to sizes that factor well.
Padding makes the array bigger. That works against saving memory. In our test
the worst case asked for many times the size of the tensor.

\subsection{Permuting more than two dimensions in place}

The closest work is ITTPD~\cite{cheng2025ittpd}. It breaks an $N$-D permutation
into block transpositions called $\rho(i,j,k)$. \textbf{That is the same
operation as our box swap. We do not claim it as ours.}

We differ in two ways.

First, ITTPD runs each step with Catanzaro's two-dimensional method. That needs
a temporary buffer of $\max(M,K)\cdot W$. It also goes over the data at least
once per step. We walk cycles directly. By Section~\ref{sec:fixedpoints} we can
go over only part of the data.

Second, ITTPD must choose its breakdown to fit a memory budget. When two
choices tie, it picks on memory rather than on time. Our memory cost is tiny
whichever breakdown we pick. So we are free to pick the cheapest one.

Where a back group of axes survives, we were faster on \resWinIttpd{} of
\resWinIttpdOf{} cases. Where it does not, ITTPD is faster than we are.
Section~\ref{sec:limits} has those numbers.

\subsection{Copying out of place}

The thing we are replacing is PyTorch's own
\texttt{.contiguous()}~\cite{pytorch_contiguous_2016,pytorch_2x}. It is fast. It
spreads perfectly over a GPU. It pays for that with a second tensor. Work on
memory pressure in large-model training~\cite{rajbhandari2020zero} and the
PyTorch communication document~\cite{pytorch_rfc_comms_2026} both treat that
cost as a real problem, not a theoretical one.

\section{Summary}
\label{sec:conclusion}

Box transpose fixes a tensor's layout in place. We tested it on \resCases{}
sampled shapes at ranks \resRankLo{} to \resRankHi{}. We also tested it on
\hardCases{} permutations picked at random. It found a working plan and gave
the right answer every time. It used \resBoxAux{} bytes of scratch space. A
copy would have used a second tensor of \resTensorMB{}\,MB.

When a permutation leaves a group of axes at the back alone, we are the fastest
in-place method we measured. We beat ITTPD, Gomez-Luna and Gustavson on
\resWinIttpd{}, \resWinGl{} and \resWinGv{} of their cases.

When the swap is also one small square block, we beat PyTorch's own copy. We
are \tierBeatsX{}$\times$ faster. The reason is simple. We skip the cells that
are already in the right place. A copy cannot.

We work out both tests from the shape, before touching any data. No number in
the cost model was measured on a GPU.

When the tests fail, our method is the slowest one we compared. That happens
when the permutation moves the innermost axis. We are about
$\hardSmallSlowdown\times$ slower there. We report it instead of routing around
it.

% Team and responsibilities. The title, authors and affiliation are set by
% \maketitle in box-transpose.tex; this file carries the division of work.
%
% A description list rather than a tabular: responsibilities are sentences, and
% a two-column table at 241pt either overflows the measure or needs its column
% widths hand-tuned every time a role changes.
\section*{Team and Responsibilities}

\begin{description}
\item[Kaveen Jayamanna (KJ).] The planner and its cost model
  (\texttt{src/\_plan.py}): the box-swap graph, the shortest-path search, and
  the closed-form fixed-point count the search is weighted by. The in-place
  cycle-following CUDA kernels (\texttt{src/\_kernels/}), including the
  vectorised and leader-based paths.

\item[Alexis Newman (AN).] Measurement and validation: the correctness suite
  that gates every result, the sweep over the shape space, the head-to-head
  against the three published in-place systems, and the adversarial sweep that
  establishes the method's limitation.
\end{description}

Both members share the write-up and the presentation.


\bibliographystyle{plain}
\bibliography{refs}

\end{document}
